\documentclass[11pt]{article}

\usepackage[a4paper,margin=27mm]{geometry}
\usepackage{amsmath,amssymb,amsthm}
\usepackage{booktabs,array}
\usepackage{xcolor}
\usepackage{listings}
\usepackage{microtype}
\usepackage{hyperref}

\hypersetup{
  colorlinks=true,
  linkcolor=blue!38!black,
  citecolor=blue!38!black,
  urlcolor=blue!48!black,
  pdftitle={An Explicit 29-AND Circuit for the AES S-box},
  pdfauthor={umizame},
  pdfsubject={An explicit twenty-nine-AND circuit for the forward AES S-box},
  pdfkeywords={AES S-box, Boolean circuit, finite fields, AND gates}
}

\definecolor{codebg}{RGB}{248,248,248}
\definecolor{codeframe}{RGB}{198,198,198}
\lstset{
  basicstyle=\ttfamily\scriptsize,
  backgroundcolor=\color{codebg},
  frame=single,
  rulecolor=\color{codeframe},
  columns=fullflexible,
  keepspaces=true,
  showstringspaces=false,
  breaklines=true,
  breakatwhitespace=false,
  aboveskip=0.8em,
  belowskip=0.8em
}

\newtheorem{theorem}{Theorem}
\newtheorem{proposition}[theorem]{Proposition}
\newtheorem{lemma}[theorem]{Lemma}

\newcommand{\F}{\mathbb F}
\newcommand{\AND}{\mathsf{AND}}
\newcommand{\XOR}{\mathsf{XOR}}
\newcommand{\gate}[1]{\mathsf A_{#1}}

\title{\bfseries An Explicit 29-AND Circuit for the AES S-box}
\author{umizame}
\date{12 July 2026; revised 27 August 2026}

\begin{document}
\maketitle

\begin{abstract}
We give an explicit straight-line program for the forward AES S-box using
twenty-nine AND gates.  The construction rests on Canright's all-normal
tower for the AES field and on the five-AND inversion of Boyar, Matthews,
and Peralta in the field of sixteen elements.  Its essential economy is a
joint calculation over the field of four elements: the three pairwise
products of three operands are obtained from three early ANDs and five late
ANDs.  Using this identity three times around the five-AND inversion gives
the schedule \(9+5+15=29\).  Affine coordinate data then realize this
schedule as a normalized XOR--AND graph and as the explicit straight-line
program.
\end{abstract}

\section{Introduction}

The nonlinear step of the AES S-box is inversion in the field of 256
elements; the remaining operations are affine over \(\F_2\).  Canright's
all-normal tower makes this inversion especially economical by resolving it
into arithmetic in successively smaller fields \cite[Sec.~2]{canright}.  Boyar,
Matthews, and Peralta subsequently gave a five-AND circuit for the inversion
which occurs at the middle level of that tower
\cite[Sec.~2.1, Fig.~2]{bmp}.  These two works
supply, respectively, the tower coordinates and the middle-level inverse
used below.

The straight-line program presented here uses the operations \(\XOR\),
\(\AND\), and \(\mathsf{NOT}\), with unrestricted fan-out.  An AND gate is
multiplication in \(\F_2\); XOR and NOT are affine and therefore require no
AND gate.  The new ingredient is a schedule for three multiplications in
\(\F_{16}\).  Their nine constituent multiplications in \(\F_4\) occur in
three triples.  In each triple the first product is needed before the
\(\F_{16}\) inverse is known, whereas the other two are needed afterwards.
A single identity computes the first product with three ANDs and completes
the other two with five more.  The nonlinear count is consequently
\[
  3\cdot(3+5)+5=29.
\]

The proof proceeds entirely in Canright's coordinates.  Section~2 recalls
the tower and the two multiplication formulae that will be used.
Section~3 proves the joint \(\F_4\) identity.  Section~4 places its three
copies around the Boyar--Matthews--Peralta inverse.  Section~5 supplies the
affine coordinates that realize the construction as a normalized XOR--AND
graph and as the explicit straight-line program.

\section{Canright's tower form}

Let
\[
 \F_{256}=\F_2[t]/(t^8+t^4+t^3+t+1),\qquad X=[t].
\]
This is the polynomial representation specified in FIPS~197
\cite[Secs.~3.2 and 4, Eqs.~(4.1) and (4.3)]{fips}.  For a byte
\(x=\sum_{i=0}^7 b_iX^i\), with \(b_i\in\F_2\), set
\[
 U_j=b_{7-j}\qquad(0\leq j<8),
\]
so that \(U_0\) is the most significant input bit.  If
\(x^{-1}=\sum_{i=0}^7v_iX^i\), with \(v_i\in\F_2\) and \(0^{-1}=0\),
FIPS~197 defines the output bits by
\begin{equation}\label{eq:fips-affine}
 o_i=v_i+v_{i+4}+v_{i+5}+v_{i+6}+v_{i+7}+c_i,
 \qquad \sum_{i=0}^7c_i2^i=\mathtt{63},
\end{equation}
where \(c_i\in\F_2\), \(\mathtt{63}\) is written in hexadecimal, and the
subscripts on \(v_i\) are read modulo eight
\cite[Sec.~5.1.1, Eqs.~(5.2)--(5.3)]{fips}.  We use
the same zero-inverse convention in every field below, and write the output
in the most-significant-first order \(S_j=o_{7-j}\).

We now take Canright's all-normal tower
\[
 \F_2\subset\F_4\subset\F_{16}\subset\F_{256},
\]
with ordered bases for the three successive quadratic extensions
\[
 [W^2,W],\qquad [Z^4,Z],\qquad [Y^{16},Y]
\]
and defining relations
\begin{equation}\label{eq:tower-relations}
 W^2+W+1=0,\qquad Z^2+Z+W^2=0,\qquad
 Y^2+Y+WZ=0.
\end{equation}
These are precisely the field and bases used in Canright's all-normal
implementation \cite[Sec.~2.1]{canright}.  The relations give
\begin{equation}\label{eq:tower-consequences}
 W^3=1,\qquad Z^4+Z=1,\qquad Z^5=W^2,\qquad
 Y^{16}+Y=1,\qquad Y^{17}=WZ.
\end{equation}
Indeed, squaring the middle relation gives
\(Z^4=Z+W^2+W=Z+1\), and multiplication by \(Z\) then gives
\(Z^5=W^2\).  At the upper level, the two roots of the quadratic in
\eqref{eq:tower-relations} are \(Y\) and \(Y+1\); the nontrivial
\(\F_{16}\)-conjugate of \(Y\) is \(Y^{16}\).  Their sum and product give
the last two identities.

Canright's change-of-basis matrix maps the tower coordinates to the AES
polynomial coordinates \cite[Sec.~4, Eq.~(10)]{canright}.  Its inverse
carries the FIPS byte to
\[
 G=G_1Y^{16}+G_0Y,\qquad G_0,G_1\in\F_{16}.
\]
The rows of this inverse map, in the bit order just fixed, are recorded in
Appendix~A.  For the formulas below, it is convenient to multiply each
\(\F_4\)-coordinate by \(W\) and put
\begin{equation}\label{eq:H-definition}
 H=WG=LY^{16}+RY,\qquad L,R\in\F_{16}.
\end{equation}
This scaling is affine, since for \(a_0,a_1\in\F_2\),
\begin{equation}\label{eq:W-scale}
 W(a_0W^2+a_1W)=(a_0+a_1)W^2+a_0W.
\end{equation}

\begin{lemma}[The upper norm]\label{lem:upper-norm}
For \(H\) as in \eqref{eq:H-definition}, define
\begin{equation}\label{eq:DQ}
 D=LR+WZ(L+R)^2,\qquad Q=D^{-1}.
\end{equation}
Then
\begin{equation}\label{eq:H-inverse}
 H^{-1}=RQY^{16}+QLY.
\end{equation}
Consequently \(W H^{-1}=G^{-1}\).
\end{lemma}

\begin{proof}
The sixteenth-power Frobenius fixes \(\F_{16}\) and exchanges
\(Y^{16}\) and \(Y\).  Thus
\[
 H^{16}=RY^{16}+LY.
\]
Using \(Y^{16}+Y=1\) and \(Y^{17}=WZ\), multiplication gives
\[
 HH^{16}=LR+WZ(L^2+R^2)=LR+WZ(L+R)^2=D.
\]
If \(H\ne0\), division by this norm proves \eqref{eq:H-inverse}.  If
\(H=0\), the basis \([Y^{16},Y]\) gives \(L=R=0\), whence \(D=Q=0\),
and the same formula holds under the declared convention.  Finally,
\(H=WG\) and \(W^3=1\), so \(WH^{-1}=G^{-1}\).
\end{proof}

The three products \(LR\), \(RQ\), and \(QL\) will be formed with the
following version of Canright's normal-basis multiplier
\cite[Sec.~2.1, Fig.~3]{canright}.

\begin{lemma}[Multiplication in \(\F_{16}\)]\label{lem:f16-product}
Let
\[
 A=A_1Z^4+A_0Z,\qquad B=B_1Z^4+B_0Z,\qquad A_i,B_i\in\F_4,
\]
and set
\begin{equation}\label{eq:f16-three-products}
 \ell=A_0B_0,\qquad h=(WA_1)(WB_1),\qquad
 s=(A_0+WA_1)(B_0+WB_1).
\end{equation}
Then
\begin{equation}\label{eq:f16-reconstruction}
 AB=(\ell+h+Ws)Z^4+(W^2h+Ws)Z.
\end{equation}
\end{lemma}

\begin{proof}
The relations in \eqref{eq:tower-consequences} give
\begin{align*}
 (Z^4)^2&=WZ^4+W^2Z,\\
 Z^2&=W^2Z^4+WZ,\\
 Z^5&=W^2Z^4+W^2Z.
\end{align*}
Hence, writing \(C=A_1B_0+A_0B_1\), direct multiplication gives
\begin{align*}
 AB={}&(WA_1B_1+W^2A_0B_0+W^2C)Z^4\\
     &+(W^2A_1B_1+WA_0B_0+W^2C)Z.
\end{align*}
Now \(h=W^2A_1B_1\) and
\(s=\ell+h+WC\).  Substitution, together with
\(1+W=W^2\), reduces the two coefficients to those in
\eqref{eq:f16-reconstruction}.
\end{proof}

\section{Three products in \texorpdfstring{\(\F_4\)}{F4}}

Write
\[
 \tau(x)=x+x^2\qquad(x\in\F_4)
\]
for the trace to \(\F_2\).  The normal basis makes both coordinates
visible as traces.

\begin{lemma}\label{lem:trace-identities}
If \(x=x_0W^2+x_1W\), with \(x_0,x_1\in\F_2\), then
\begin{equation}\label{eq:trace-coordinates}
 \tau(x)=x_0+x_1,\qquad \tau(Wx)=x_1,\qquad
 \tau(W^2x)=x_0.
\end{equation}
For all \(x,y\in\F_4\),
\begin{align}
 \tau(x)\tau(Wy)+\tau(Wx)\tau(y)&=\tau(xy^2),
 \label{eq:trace-first}\\
 \tau(x)\tau(y)+\tau(Wx)\tau(Wy)&=\tau(Wxy).
 \label{eq:trace-second}
\end{align}
\end{lemma}

\begin{proof}
Squaring \(x_0W^2+x_1W\) exchanges its two normal-basis coordinates,
which proves \eqref{eq:trace-coordinates}.  If
\(y=y_0W^2+y_1W\), with \(y_0,y_1\in\F_2\), the left side of
\eqref{eq:trace-first} is
\(x_0y_1+x_1y_0\), while the left side of
\eqref{eq:trace-second} is
\(x_0y_0+x_0y_1+x_1y_0\).  Multiplication with \(W^2+W=1\) shows that
these are respectively \(\tau(xy^2)\) and \(\tau(Wxy)\).
\end{proof}

Let \(a,b,c\in\F_4\).  We first recover \(ab\) with three ANDs that do
not depend on \(c\); five later ANDs will then recover \(bc\) and \(ca\)
jointly.  The three early products are
\begin{equation}\label{eq:early-products}
 p_1=\tau(a)\tau(b),\qquad
 p_2=\tau(Wa)\tau(Wb),\qquad
 p_3=\tau(W^2a)\tau(W^2b).
\end{equation}
Since \(\tau(W^2x)=\tau(x)+\tau(Wx)\), equations
\eqref{eq:trace-first} and \eqref{eq:trace-second} give
\begin{equation}\label{eq:sigma-ab}
 \sigma:=p_1+p_2+p_3=\tau(ab^2),\qquad
 ab=(p_1+p_3)W^2+(p_1+p_2)W.
\end{equation}
For the second formula, the coefficient of \(W\) is
\(\tau(Wab)=p_1+p_2\); applying \eqref{eq:trace-second} after cycling
both arguments by \(W^2\) gives the coefficient of \(W^2\) as
\(\tau(W^2ab)=p_1+p_3\).

For distinct \(r,t\in\{1,W,W^2\}\), put
\begin{equation}\label{eq:cell}
 e_{r,t}=
 \bigl(\tau(ra)+\tau(tb)+\tau(Wc)\bigr)
 \bigl(\tau(Wra)+\tau(Wtb)+\tau(c)\bigr).
\end{equation}
Among the six off-diagonal cells \(e_{r,t}\), the five \(e\)-terms used
below form the union of the rows \(r=W,W^2\) and the columns \(t=1,W^2\).

\begin{proposition}[The five late products]\label{prop:f4-joint}
In addition to \(p_1,p_2,p_3\), define
\begin{equation}\label{eq:late-products}
\begin{aligned}
 p_4&=e_{1,W^2},&
 p_5&=e_{W,W^2},\\
 p_6&=e_{W^2,1}+\tau(a+Wb+c),&
 p_7&=e_{W,1},\\
 p_8&=e_{W^2,W}+\tau(a+W^2b+c).
\end{aligned}
\end{equation}
Each \(p_i\) in \eqref{eq:late-products} is the output of one AND gate.
Moreover,
\begin{align}
 \tau(W^2bc)&=\sigma+p_5+p_7+\tau(b),
 \label{eq:bc-first}\\
 \tau(bc)&=\sigma+p_6+p_8+\tau(W^2b),
 \label{eq:bc-second}\\
 \tau(Wca)&=\sigma+p_6+p_7+\tau(a+Wb+c)+\tau(W^2a),
 \label{eq:ca-first}\\
 \tau(ca)&=\sigma+p_4+p_5+\tau(Wa).
 \label{eq:ca-second}
\end{align}
Thus \(ab\), \(bc\), and \(ca\) are obtained from eight AND gates,
the first three of which do not depend on \(c\).
\end{proposition}

\begin{proof}
We first compute the sum of two cells in a common row and in a common
column; the stated formulas then follow after accounting for the two affine
shears.
Write \(\tau_r(x)=\tau(rx)\) and
\(\eta(x)=\tau(x)\tau(Wx)\).  The three bits
\(\tau_1(x),\tau_W(x),\tau_{W^2}(x)\) sum to zero.  If
\(\{r,t,u\}=\{1,W,W^2\}\), expansion of the two cells in a fixed row
gives
\begin{equation}\label{eq:row-expansion}
\begin{aligned}
 e_{r,t}+e_{r,u}={}&
  \tau_r(a)\tau_{Wr}(b)+\tau_{Wr}(a)\tau_r(b)\\
 &+\tau_W(c)\tau_{Wr}(b)+\tau_1(c)\tau_r(b)\\
 &+\tau_t(b)\tau_{Wt}(b)+\tau_u(b)\tau_{Wu}(b).
\end{aligned}
\end{equation}
The first line is \(\tau(ab^2)=\sigma\) by
\eqref{eq:trace-first}, since \(r^3=1\).  The second is
\(\tau(Wrbc)\) by \eqref{eq:trace-second}.  The last line is
\(\eta(tb)+\eta(ub)\).  Writing \(x=x_0W^2+x_1W\), with
\(x_0,x_1\in\F_2\), gives
\[
\eta(x)=x_1(x_0+x_1)=x_0x_1+x_1
        =x^3+\tau(W^2x).
\]
Here \(x^3\) is the norm from \(\F_4\) to \(\F_2\).
Since \(t^3=u^3=1\), the two norm terms cancel; since \(t+u=r\),
the remaining trace is \(\tau(W^2rb)\).  Thus
\begin{equation}\label{eq:row-identity}
 e_{r,t}+e_{r,u}
 =\sigma+\tau(Wrbc)+\tau(W^2rb).
\end{equation}

The same argument down a column gives
\begin{align*}
 e_{r,t}+e_{u,t}={}&
  \tau_t(b)\tau_{Wt}(a)+\tau_{Wt}(b)\tau_t(a)\\
 &+\tau_W(c)\tau_{Wt}(a)+\tau_1(c)\tau_t(a)\\
 &+\tau_r(a)\tau_{Wr}(a)+\tau_u(a)\tau_{Wu}(a),
\end{align*}
and therefore
\begin{equation}\label{eq:column-identity}
 e_{r,t}+e_{u,t}
 =\sigma+\tau(Wtca)+\tau(W^2ta).
\end{equation}
Here the first line is again \(\sigma\) by
\eqref{eq:trace-first}; explicitly,
\(\tau(ba^2)=\tau((ba^2)^2)=\tau(ab^2)\), by Frobenius invariance of
the trace.  The second line follows from \eqref{eq:trace-second}.  The last
line is
\(\eta(ra)+\eta(ua)\); cancellation of the two norms and
\(r+u=t\) leave \(\tau(W^2ta)\).

Equations~\eqref{eq:row-identity} and \eqref{eq:column-identity} organize
the reconstruction: row sums give traces of \(bc\), whereas column sums
give traces of \(ca\).

It remains to account for the two sheared cells.  Write the two factors of
\(e_{W^2,1}\) as \(f\) and \(g\).  The second is
\(g=\tau(a+Wb+c)\); hence
\[
 p_6=fg+g=(f+g)g,
\]
because \(g^2=g\) in \(\F_2\).  Thus \(p_6\) is one AND, with factors
\(f+g\) and \(g\).  The second factor of \(e_{W^2,W}\) is
\(\tau(a+W^2b+c)\), and the same argument proves the assertion for
\(p_8\).

Taking \(r=W\) and columns \(W^2,1\) in
\eqref{eq:row-identity} gives \eqref{eq:bc-first}.  Taking \(r=W^2\)
and columns \(1,W\), and observing that the two shear terms sum to
\(\tau(b)\), while
\(\tau(Wb)+\tau(b)=\tau(W^2b)\), gives \eqref{eq:bc-second}.  Taking
column \(1\) and rows
\(W^2,W\) in \eqref{eq:column-identity} gives
\eqref{eq:ca-first}; taking column \(W^2\) and rows \(1,W\) gives
\eqref{eq:ca-second}.

Finally, \eqref{eq:trace-coordinates} gives the explicit reconstructions
\[
 \begin{split}
 bc&=\tau(W^2bc)W^2+
       \bigl(\tau(bc)+\tau(W^2bc)\bigr)W,\\
 ca&=\bigl(\tau(ca)+\tau(Wca)\bigr)W^2+\tau(Wca)W.
 \end{split}
\]
Thus the four displayed traces recover \(bc\) and \(ca\), while
\eqref{eq:sigma-ab} has already recovered \(ab\).
\end{proof}

\section{The twenty-nine AND gates}

Write
\[
 L=L_1Z^4+L_0Z,\qquad
 R=R_1Z^4+R_0Z,\qquad
 Q=Q_1Z^4+Q_0Z,\qquad L_i,R_i,Q_i\in\F_4.
\]
Apply Proposition~\ref{prop:f4-joint} to the three triples
\begin{equation}\label{eq:three-rows}
 \begin{array}{ccc}
 a&b&c\\
 \hline
 L_0&R_0&Q_0\\
 WL_1&WR_1&WQ_1\\
 L_0+WL_1&R_0+WR_1&Q_0+WQ_1.
 \end{array}
\end{equation}
In the three rows, the products \(ab\) are respectively the quantities
\(\ell,h,s\) of Lemma~\ref{lem:f16-product} for \(LR\).  The products
\(bc\) are the same three quantities for \(RQ\), and the products
\(ca\) are those for \(QL\).  Thus the first three ANDs in each row
make \(LR\) available; the final five, once \(Q\) is known, complete
\(RQ\) and \(QL\).

It remains to obtain \(Q=D^{-1}\) between these two stages.  We use the
five-product inversion of Boyar, Matthews, and Peralta, written in
Canright's coordinates.

\begin{proposition}[The middle inverse]\label{prop:middle-inverse}
Write
\begin{equation}\label{eq:D-coordinates}
 D=(d_0W^2+d_1W)Z^4+(d_2W^2+d_3W)Z.
\end{equation}
Here \(d_0,\ldots,d_3\in\F_2\).
Define, in order,
\begin{equation}\label{eq:bmp-products}
\begin{aligned}
 m_1&=(d_0+d_1)(d_2+d_3),\\
 m_2&=d_1(d_2+m_1),\\
 m_3&=(d_0+m_1)d_3,\\
 m_4&=d_2(m_1+m_3),\\
 m_5&=(d_0+m_2)(d_2+m_1+m_4).
\end{aligned}
\end{equation}
Also define
\begin{equation}\label{eq:Q-coordinates}
 \begin{split}
 q_0&=d_2+d_3+m_4,\qquad
 q_1=d_3+m_3+m_4,\\
 q_2&=d_0+d_1+m_2+m_5,\qquad
 q_3=d_1+m_5.
 \end{split}
\end{equation}
Then the five products in \eqref{eq:bmp-products} give
\[
 Q=(q_0W^2+q_1W)Z^4+(q_2W^2+q_3W)Z=D^{-1}.
\]
\end{proposition}

\begin{proof}
Boyar, Matthews, and Peralta represent \(\Delta\) using bits
\(x_1,\ldots,x_4\in\F_2\) as
\[
 \Delta=(x_1W+x_2W^2)Z^2+(x_3W+x_4W^2)Z^8
\]
and give a five-AND circuit for \(\Delta^{-1}\)
\cite[Sec.~2.1, Fig.~2]{bmp}.  Set
\begin{equation}\label{eq:bmp-change}
 (x_1,x_2,x_3,x_4)=(d_0+d_1,d_0,d_2+d_3,d_2).
\end{equation}
Since \(Z^{16}=Z\), squaring \(\Delta\) and multiplying by \(W^2\)
gives \(D=W^2\Delta^2\).  Indeed,
\(W^2((d_0+d_1)W+d_0W^2)^2=d_0W^2+d_1W\), and the
same identity holds for the lower pair.
Taking inverses and using \(W^3=1\) then gives
\begin{equation}\label{eq:D-Delta}
 D=W^2\Delta^2,\qquad Q=D^{-1}=W(\Delta^{-1})^2.
\end{equation}

Under \eqref{eq:bmp-change}, the five ANDs in that circuit are exactly
\(m_1,\ldots,m_5\) in
\eqref{eq:bmp-products}.  Its four output coordinates are
\begin{align*}
 y_1&=d_3+m_3+m_4,& y_2&=d_2+m_3,\\
 y_3&=d_1+m_5,& y_4&=d_0+m_2,
\end{align*}
where
\[
 \Delta^{-1}=(y_1W+y_2W^2)Z^2+(y_3W+y_4W^2)Z^8.
\]
For \(\alpha=y_1W+y_2W^2\), one has
\[
 W\alpha^2=(y_1+y_2)W^2+y_1W.
\]
Applying this to both pairs in \eqref{eq:D-Delta} yields precisely the
four expressions in \eqref{eq:Q-coordinates}.  The five-product circuit
returns zero on zero, so the identity also respects the convention fixed
above.
\end{proof}

We may now read off the entire schedule.  The three early gates in each row of
\eqref{eq:three-rows} use only \(L\) and \(R\), and nine ANDs therefore
produce \(LR\).  Squaring and multiplication by the constant \(WZ\) are
linear, so \(D=LR+WZ(L+R)^2\) is then available without another AND.
Proposition~\ref{prop:middle-inverse} produces \(Q\) with five ANDs.  The
five late gates in each row of \eqref{eq:three-rows} then produce \(RQ\)
and \(QL\) with fifteen further ANDs.  Lemma~\ref{lem:upper-norm} gives
\(H^{-1}\).  Multiplication by \(W\), the combined linear map from the tower
coordinates to the linear part of the FIPS output, and addition of the FIPS
constant are all affine.  This schedule is acyclic, and its nonlinear tally is
\[
 9+5+15=29.
\]

\section{The explicit circuit}\label{sec:identification}

To separate the nonlinear schedule from a particular affine XOR network,
Appendix~B records the construction as a normalized XOR--AND graph (XAG).
This XAG uses the
ordered affine basis
\[
 (1,U_0,\ldots,U_7,\gate{1},\ldots,\gate{29}).
\]
For an eight-bit linear mask \(\mu=\sum_{j=0}^7\mu_j2^j\), which is
written in hexadecimal when displayed, define
\[
 \lambda_\mu=\sum_{j=0}^7\mu_jU_j;
\]
thus bit \(j\) selects \(U_j\), and the least significant bit selects
\(U_0\).  By contrast, the XAG uses affine masks with an additional
least-significant constant bit.  In an AND row for \(\gate{i}\), bit zero
selects the constant, bits one through eight select \(U_0,\ldots,U_7\),
and bit \(8+j\) selects \(\gate{j}\) for \(1\leq j<i\).  The product of
the two affine forms is \(\gate{i}\).  An output row is an affine mask in
the complete basis.

Appendix~A gives all coordinate data needed to express the preceding
construction in these masks, together with the gate ranges for the three
rows of \eqref{eq:three-rows}.  The five products
\(m_1,\ldots,m_5\) are \(\gate{10},\ldots,\gate{14}\).

To illustrate the substitution, the lower row has
\[
 L_0=\lambda_{\mathtt{41}}W^2+\lambda_{\mathtt{c6}}W,\qquad
 R_0=\lambda_{\mathtt{74}}W^2+\lambda_{\mathtt{86}}W.
\]
Encoding these forms in the affine-mask convention above gives the three
early factor pairs, in the orientation used in Appendix~B,
\[
 (\mathtt{1e4},\mathtt{10e}),\qquad
 (\mathtt{10c},\mathtt{18c}),\qquad
 (\mathtt{e8},\mathtt{82}).
\]
These are the first three rows of Appendix~B.
The same substitution gives \(\gate{4}\) through \(\gate{9}\).
The four masks of \(D\) in Appendix~A give
\[
 (d_0+d_1,d_2+d_3)=(\mathtt{1ecf6},\mathtt{1d06a}),
\]
the factor pair for \(\gate{10}\); successive substitution in
\eqref{eq:bmp-products} gives \(\gate{11}\) through \(\gate{14}\).

Substituting the four masks of \(Q\) into \eqref{eq:late-products} in each
of the three rows gives the fifteen late rows in the order tabulated in
Appendix~A.  Finally, Lemma~\ref{lem:f16-product}
reconstructs \(RQ\) and \(QL\), and the affine rows in Appendix~A express
their coordinates as the eight output masks printed in Appendix~B.
Thus every row of Appendix~B is obtained from the displayed formulae by
successive XOR substitution.  A concrete affine realization of this XAG is
the straight-line program supplied in the repository:
\begin{center}
 \href{https://github.com/umizame/S-box_29-AND/blob/master/circuits/aes-sbox-fwd-g228-a29-d35-ad6.slp}{\texttt{circuits/aes-sbox-fwd-\allowbreak g228-a29-d35-ad6.slp}}.
\end{center}

Scanning its assignments in order expresses every wire as an affine form in
the inputs, the constant, and the outputs of preceding AND gates.  XOR and
NOT update the affine mask; each AND assignment records its two factor masks
and introduces the next \(\gate{i}\).  Consequently the program's AND
assignments give, in order, the twenty-nine factor rows in Appendix~B, and
its final affine forms give the eight output rows.

\begin{theorem}\label{thm:main}
The straight-line program above computes the forward AES S-box specified in
FIPS~197 and contains twenty-nine AND gates.
\end{theorem}

\begin{proof}
The inverse of Canright's change-of-basis matrix sends the AES polynomial
coordinates to \(G\), and \eqref{eq:H-definition} gives \(H=WG\).
Lemma~\ref{lem:f16-product},
Proposition~\ref{prop:f4-joint}, and
Proposition~\ref{prop:middle-inverse} construct \(RQ\) and \(QL\) with
twenty-nine AND gates.  Lemma~\ref{lem:upper-norm} then gives
\(WH^{-1}=G^{-1}\), including at zero.  The affine coordinates in
Appendix~A apply multiplication by \(W\), Canright's tower-to-polynomial
map, and the linear part of the FIPS affine map, before adding the constant in
\eqref{eq:fips-affine}.  They therefore give the FIPS output in the declared
bit order and the normalized XAG in Appendix~B.  The straight-line program
above is a concrete affine realization of that XAG, so it computes the same
output; its twenty-nine AND assignments are precisely the twenty-nine AND
gates of the construction.
\end{proof}

\section*{Acknowledgement}

GPT-5.6 Pro assisted with algebraic checking and with the preparation of the
exposition.

\appendix

\section{Affine coordinates}

We use both hexadecimal-mask conventions of
Section~\ref{sec:identification}.  On the ordered linear input basis
\((U_0,\ldots,U_7)\), bit \(j\) selects \(U_j\).  On an ordered affine
basis \((1,V_0,\ldots,V_{n-1})\), bit zero selects the constant and bit
\(j+1\) selects \(V_j\).  All masks below are written in hexadecimal.

The rows of the inverse map from AES polynomial coordinates to Canright's
tower coordinates, and the corresponding rows after the scaling \(H=WG\),
are as follows \cite[Sec.~4, Eq.~(10)]{canright}.  The
coordinate order is the tensor order induced by
\([Y^{16},Y]\), \([Z^4,Z]\), and \([W^2,W]\).
\begin{center}
\begin{tabular}{@{}ccc@{}}
\toprule
coordinate&\(G\)&\(H=WG\)\\
\midrule
\(Y^{16}Z^4W^2\)&\(\mathtt{e7}\)&\(\mathtt{69}\)\\
\(Y^{16}Z^4W\)&\(\mathtt{8e}\)&\(\mathtt{e7}\)\\
\(Y^{16}ZW^2\)&\(\mathtt{c6}\)&\(\mathtt{41}\)\\
\(Y^{16}ZW\)&\(\mathtt{87}\)&\(\mathtt{c6}\)\\
\(YZ^4W^2\)&\(\mathtt{d9}\)&\(\mathtt{59}\)\\
\(YZ^4W\)&\(\mathtt{80}\)&\(\mathtt{d9}\)\\
\(YZW^2\)&\(\mathtt{86}\)&\(\mathtt{74}\)\\
\(YZW\)&\(\mathtt{f2}\)&\(\mathtt{86}\)\\
\bottomrule
\end{tabular}
\end{center}
In particular,
\begin{align*}
 L_1&=\lambda_{\mathtt{69}}W^2+\lambda_{\mathtt{e7}}W,&
 L_0&=\lambda_{\mathtt{41}}W^2+\lambda_{\mathtt{c6}}W,\\
 R_1&=\lambda_{\mathtt{59}}W^2+\lambda_{\mathtt{d9}}W,&
 R_0&=\lambda_{\mathtt{74}}W^2+\lambda_{\mathtt{86}}W.
\end{align*}
The three rows of \eqref{eq:three-rows}, with their gate ranges, are
\begin{center}
\begin{tabular}{@{}ccccc@{}}
\toprule
row&\(a\)&\(b\)&\(c\)&early; late\\
\midrule
lower&\(L_0\)&\(R_0\)&\(Q_0\)&\(\gate{1}\text{--}\gate{3}\); \(\gate{15}\text{--}\gate{19}\)\\
scaled upper&\(WL_1\)&\(WR_1\)&\(WQ_1\)&\(\gate{4}\text{--}\gate{6}\); \(\gate{20}\text{--}\gate{24}\)\\
sum&\(L_0+WL_1\)&\(R_0+WR_1\)&\(Q_0+WQ_1\)&\(\gate{7}\text{--}\gate{9}\); \(\gate{25}\text{--}\gate{29}\)\\
\bottomrule
\end{tabular}
\end{center}
The input masks for \((a,b)\) in these rows, as pairs of
\([W^2,W]\)-coordinates, are
\[
 (\mathtt{41},\mathtt{c6}),\ (\mathtt{74},\mathtt{86});\qquad
 (\mathtt{8e},\mathtt{69}),\ (\mathtt{80},\mathtt{59});\qquad
 (\mathtt{cf},\mathtt{af}),\ (\mathtt{f4},\mathtt{df}).
\]
Within every late range, the gates are \(p_4,p_5,p_6,p_7,p_8\) in that
order; \(p_6\) and \(p_8\) use the sheared factors displayed in the proof
of Proposition~\ref{prop:f4-joint}.

On the affine basis
\((1,U_0,\ldots,U_7,\gate{1},\ldots,\gate{9})\), the coordinate masks of
\(D\) in \eqref{eq:D-coordinates} are
\begin{center}
\begin{tabular}{@{}cc@{}}
\toprule
coordinate&mask\\
\midrule
\(d_0\)&\(\mathtt{35afc}\)\\
\(d_1\)&\(\mathtt{2b60a}\)\\
\(d_2\)&\(\mathtt{330ea}\)\\
\(d_3\)&\(\mathtt{2e080}\)\\
\bottomrule
\end{tabular}
\end{center}
On the basis enlarged through \(\gate{14}\), the four coordinate masks of
\(Q\) are
\begin{center}
\begin{tabular}{@{}cc@{}}
\toprule
coordinate&mask\\
\midrule
\(q_0\)&\(\mathtt{21d06a}\)\\
\(q_1\)&\(\mathtt{32e080}\)\\
\(q_2\)&\(\mathtt{49ecf6}\)\\
\(q_3\)&\(\mathtt{42b60a}\)\\
\bottomrule
\end{tabular}
\end{center}

For the output coordinates, write
\begin{align*}
 RQ&=(T_0W^2+T_1W)Z^4+(T_2W^2+T_3W)Z,\\
 QL&=(T_4W^2+T_5W)Z^4+(T_6W^2+T_7W)Z.
\end{align*}
Here \(T_0,\ldots,T_7\in\F_2\).
In the same most-significant-first convention, composing Canright's
tower-to-polynomial map \cite[Sec.~4, Eq.~(10)]{canright} with the linear
part of the FIPS affine map
\cite[Sec.~5.1.1, Eq.~(5.4)]{fips} gives the rows
\[
 (\mathtt{14},\mathtt{11},\mathtt{82},\mathtt{15},
   \mathtt{1f},\mathtt{b6},\mathtt{4c},\mathtt{4a}).
\]
Composing multiplication by \(W\) with this linear map and then adding the
constant \(\mathtt{63}\) gives the following affine masks on
\((1,T_0,\ldots,T_7)\):
\begin{center}
\begin{tabular}{@{}cccccccc@{}}
\toprule
\(S_0\)&\(S_1\)&\(S_2\)&\(S_3\)&\(S_4\)&\(S_5\)&\(S_6\)&\(S_7\)\\
\midrule
\(\mathtt{078}\)&\(\mathtt{067}\)&\(\mathtt{083}\)&\(\mathtt{07e}\)&
\(\mathtt{074}\)&\(\mathtt{0da}\)&\(\mathtt{191}\)&\(\mathtt{18b}\)\\
\bottomrule
\end{tabular}
\end{center}
After expressing the \(T_i\) on the complete XAG basis, these become
\begin{center}
\begin{tabular}{@{}cc@{\qquad}cc@{}}
\toprule
output&mask&output&mask\\
\midrule
\(S_0\)&\(\mathtt{2491ba35ce}\)&\(S_1\)&\(\mathtt{257bba3b95}\)\\
\(S_2\)&\(\mathtt{2e6532ee01}\)&\(S_3\)&\(\mathtt{19dbba4b98}\)\\
\(S_4\)&\(\mathtt{189eba357c}\)&\(S_5\)&\(\mathtt{c0332e152}\)\\
\(S_6\)&\(\mathtt{db00b2bf9}\)&\(S_7\)&\(\mathtt{19b508a44b}\)\\
\bottomrule
\end{tabular}
\end{center}

\section{The normalized XAG}

\lstinputlisting[caption={The normalized XAG
\texttt{certificates/aes29.xag}.},label={lst:xag}]{certificates/aes29.xag}

\begin{thebibliography}{9}

\bibitem{canright}
D. Canright,
``A Very Compact S-box for AES,''
in \emph{Cryptographic Hardware and Embedded Systems -- CHES 2005},
Lecture Notes in Computer Science, vol.~3659, pp.~441--455,
Springer, 2005.
\href{https://doi.org/10.1007/11545262_32}{doi:10.1007/11545262\_32}.

\bibitem{bmp}
J. Boyar, P. Matthews, and R. Peralta,
``Logic Minimization Techniques with Applications to Cryptology,''
\emph{Journal of Cryptology}, vol.~26, pp.~280--312, 2013.
\href{https://doi.org/10.1007/s00145-012-9124-7}{doi:10.1007/s00145-012-9124-7}.

\bibitem{fips}
National Institute of Standards and Technology,
\emph{Advanced Encryption Standard (AES)},
NIST FIPS 197-upd1, originally published 2001, updated 9 May 2023.
\href{https://doi.org/10.6028/NIST.FIPS.197-upd1}{doi:10.6028/NIST.FIPS.197-upd1}.

\end{thebibliography}

\end{document}
